% ══════════════════════════════════════════════════════════════════════════
\chapter{Experimento C: incidencia con y sin pseudo-Brewster}
\label{cap:resC}

\paragraph{Hipótesis (§2.1).} Cerca del pseudo-Brewster se \emph{reduce} el
retorno de la interfaz en TM (no aumenta el eco del blanco), lo que mejoraría
la relación blanco--superficie si esa reducción es robusta a pérdidas,
rugosidad, errores angulares y de atitud.

\paragraph{Método.} Se comparan dos familias de trayectoria del Tx (F1, hélice; F2,
incidencia constante), descritas abajo, en varias incidencias, suelos, rugosidades,
polarizaciones y errores de apunte (Tabla~\ref{tab:C_var}). Canales: mono, bi a 90°, 135°
y 180° (Ref1), y fusión. El SCR se mide en el vóxel del blanco (punto de
$1\,\mathrm{m}^2$ a 2\,m de profundidad), con la superficie congelada y compartida por los
dos canales (3 realizaciones), ponderación Hann y ventana de recepción
(Sección~\ref{sec:ponderacion}). CSV:
\texttt{expC\_enlace\_reflectancia.csv}, \texttt{expC\_enlace\_error\_roll.csv},
\texttt{expC\_scr\_imagen.csv}, \texttt{expC\_psf\_vs\_incidencia.csv}; versión
con espectro rectangular en \texttt{expC\_rect/}.

\paragraph{Geometría de F1 y F2.} Las dos familias se construyen punto a punto a partir de
la hélice base (\texttt{geometry.familia\_C}). Cada pulso conserva su acimut y su distancia al
centro de la escena, que en la hélice óptima es casi constante ($\approx189$\,m, porque su línea
de base es perpendicular a la línea de vista). Sólo cambia la incidencia $\theta$, medida en el
centro de la superficie $(0,0,0)$:
\begin{itemize}[itemsep=1pt]
  \item \textbf{F1:} cada punto se gira en su plano vertical $\theta_c-\langle\theta_0\rangle$. La
  trayectoria sigue siendo una hélice y conserva su barrido propio de incidencia
  ($\theta_c\pm7^{\circ}$).
  \item \textbf{F2:} cada punto se coloca exactamente en $\theta_c$. Como la distancia es casi
  constante, F2 es una órbita de dos vueltas a altura y radio constantes; en
  $\theta_{pB}=63.45^{\circ}$, 84.6\,m de altura y 169\,m de radio.
\end{itemize}
Un \emph{cono de incidencia constante} $\theta$ es el conjunto de puntos que, vistos desde
el centro de la superficie, forman el mismo ángulo $\theta$ con la vertical: un cono con vértice
en $(0,0,0)$ y semiángulo $\theta$ (Fig.~\ref{fig:C_3d}). La hélice base y F1 lo atraviesan
(su incidencia cambia en cada pulso); F2 vive sobre él.

\begin{figure}[H]
\centering
\includegraphics[width=\textwidth]{C_trayectorias_3d.pdf}
\caption{Trayectorias completas del Tx en 3D (\texttt{exp\_C\_fig\_trayectorias3d.py}). (a)
Hélice base, entre los conos de $46.3^{\circ}$ y $60.5^{\circ}$. (b) F1 en $\theta_{pB}$: la
misma hélice girada hacia abajo, que cruza el cono de $\theta_{pB}$ (rojo) a mitad de la
apertura. (c) F2 en $\theta_{pB}$: todos los pulsos sobre el cono de $63.45^{\circ}$, es decir,
un círculo de 169\,m de radio a 85\,m de altura recorrido dos veces.}
\label{fig:C_3d}
\end{figure}

El Rx sigue a su Tx a la misma altura y el mismo radio, desfasado $\Delta\phi$, como en el
Exp.~B. En F2 ve, por tanto, la misma incidencia $\theta_c$.

La reflexión especular no ocurre en $\theta_c$, sino en el punto especular del par Tx--Rx, que
está a mitad de camino entre sus nadires (Fig.~\ref{fig:C_geo}\,c). Su incidencia es
\begin{equation}
  \theta_{esp}=\arctan\!\big(\tan\theta_c\,\sin(\Delta\phi/2)\big),
  \label{eq:C_esp}
\end{equation}
y está a $\rho\,|\cos(\Delta\phi/2)|$ del centro. En el canal mono la especular es el reflejo de
nadir ($\theta_{esp}=0$); a $90^{\circ}$, con $\theta_c=\theta_{pB}$, $\theta_{esp}=54.8^{\circ}$; a
$135^{\circ}$, $61.6^{\circ}$. Sólo en Ref1 ($180^{\circ}$) coincide con $\theta_c$
(Fig.~\ref{fig:C_geo}\,d). \textbf{Por eso, elegir $\theta_c=\theta_{pB}$ sólo pone la
reflexión especular en Brewster en Ref1.} En las demás configuraciones, lo que se suprime en
$\theta_c$ es el retorno difuso. El modelo calcula la especular en el punto especular real de
cada pulso (\texttt{scene.specular\_echo}, plano infinito), no en $\theta_c$.

\begin{figure}[H]
\centering
\includegraphics[width=\textwidth]{C_geometria.pdf}
\caption{Geometría del Exp.~C (\texttt{exp\_C\_fig\_geometria.py}). (a) Plano vertical
$(\rho,z)$: hélice base, F1 en $\theta_{pB}$ (segmento de $\pm7^{\circ}$) y F2 (un punto por
$\theta_c$, con los errores $\theta_{pB}+\delta$ y las incidencias fuera de Brewster). (b)
Incidencia del Tx en cada pulso. (c) Planta de F2 en $\theta_{pB}$, con el Rx a 90°, 135° y
180° y el punto especular de cada par. (d) Incidencia en el punto especular frente a
$\Delta\phi$: ecuación (\ref{eq:C_esp}) y valores de la simulación.}
\label{fig:C_geo}
\end{figure}

\begin{table}[H]
\centering\small
\caption{Variables del Exp.~C y cómo entran en la simulación (\texttt{params/expC.json},
\texttt{exp\_C\_brewster.py}).}
\label{tab:C_var}
\begin{tabular}{@{}L{2.5cm}L{3.3cm}L{8.6cm}@{}}
\toprule
\textbf{Variable} & \textbf{Valores} & \textbf{Qué cambia en la simulación} \\
\midrule
Familia & F1, F2 & Forma de la trayectoria del Tx (y del Rx): hélice con barrido de $\pm7^{\circ}$
o cono de incidencia constante \\
$\theta_c$ & 35°, 45°, 53.4° (base), 75°; $\theta_{pB}+\delta$ & Incidencia central del Tx en el
centro de la superficie; fija la altura y el radio de la órbita \\
$\theta_{pB}$ & S0 63.43°, S1 63.45°, S2 68.66°, S3 74.05° & Mínimo de $|r_{TM}(\theta)|^2$ del suelo;
sólo sirve para elegir $\theta_c$ \\
$\delta$ (error angular) & 0, $\pm0.5$, $\pm1$, $\pm2$, $\pm5$° & Desplaza toda la trayectoria:
$\theta_c=\theta_{pB}+\delta$ (error sistemático de apunte o de altura) \\
$\Delta\phi$ & 0 (mono), 90°, 135°, 180° & Rx en la misma órbita desfasado $\Delta\phi$; fija
$\theta_{esp}$ (ec.~\ref{eq:C_esp}) y la posición del punto especular \\
Polarización & TM, TE & Coeficiente de Fresnel en la especular, núcleo SPM ($\approx$VV o HH) en la
difusa y transmitancia del blanco; Tx y Rx co-polares \\
Roll $\gamma$ & 0, 1, 2, 5° & Antena girada alrededor de la línea de vista: pesos
$\cos^2\gamma$ (TM) y $\sin^2\gamma$ (TE); sólo F2 en $\theta_{pB}$, S1, TM \\
Suelo & S0--S3 (S0 4, S1 3.99$-j$0.40, S2 6.5$-j$0.9, S3 12$-j$2.5) & Fresnel, $\sigma^0$ SPM,
transmitancia y $\theta_{pB}$. S1 completo; S0, S2 y S3 en $\theta_{pB}$, $\theta_{pB}\pm2^{\circ}$ y 53.4° \\
Rugosidad $s$ & 0, 1, 2, 4\,cm ($\ell=20$\,cm) & Especular $\times\,e^{-2k_0^2s^2\cos^2\theta_{esp}}$
(componente coherente); difusa SPM con amplitud $\propto s$ sobre un parche de $\pm12$\,m
(facetas de 15\,cm, 3 realizaciones) \\
SCR & enlace e imagen & Enlace: potencias $P_{blanco}/(P_{esp}+P_{dif})$. Imagen:
$|I_{blanco}|^2/\langle|I_{sup}|^2\rangle$ en el vóxel del blanco tras la back-projection
(ventana Hann y ventana de recepción) \\
\bottomrule
\end{tabular}
\end{table}

\section{Enlace: reflectancia y su robustez}

\begin{figure}[H]
\centering
\includegraphics[width=\textwidth]{C_reflectancia.pdf}
\caption{Reflectancia TM y TE de los suelos (izq.) y reflectancia TM
efectiva frente al error de incidencia y de atitud para S0 (sin pérdidas) y
S1 ($\varepsilon_r=3.99-j0.40$).}
\label{fig:C_refl}
\end{figure}

\begin{table}[H]
\centering\small
\caption{Reflectancia TM efectiva [dB] en $\theta_{pB}+\delta$ (\texttt{expC\_enlace\_error\_roll.csv}).
Roll 0°: media de $\pm\delta$ (simétrica a $<1$\,dB); roll 5°: pares $+\delta$\,/\,$-\delta$.
Última columna: incidencia base 53.4°. Menor es mejor (más supresión). Color según la supresión
respecto de la incidencia base (en el peor de los dos signos de $\delta$):
\colorbox{mejor}{$\ge10$\,dB}, \colorbox{regular}{5--10\,dB}, \colorbox{peor}{$<5$\,dB}.}
\label{tab:C_err}
\resizebox{\textwidth}{!}{\begin{tabular}{@{}lrrrrrr|r@{}}
\toprule
Suelo / roll & $\delta=0$ & $\pm0.5^{\circ}$ & $\pm1^{\circ}$ & $\pm2^{\circ}$ & $\pm5^{\circ}$ & & fuera (53.4°) \\
\midrule
S0, roll 0° & \cellcolor{mejor}$-189.7$ & \cellcolor{mejor}$-41.7$ & \cellcolor{mejor}$-35.7$ & \cellcolor{mejor}$-29.7$ & \cellcolor{peor}$-21.7$ & & $-17.5$ \\
S0, roll 5° & \cellcolor{mejor}$-46.8$ & \cellcolor{mejor}$-37.8$/$-49.1$ & \cellcolor{mejor}$-33.4$/$-38.9$ & \cellcolor{mejor}$-28.2$/$-31.5$ & \cellcolor{peor}$-20.3$/$-23.3$ & & --- \\
S1, roll 0° & \cellcolor{mejor}$-34.5$ & \cellcolor{mejor}$-33.8$ & \cellcolor{mejor}$-32.1$ & \cellcolor{mejor}$-28.5$ & \cellcolor{peor}$-21.5$ & & $-17.5$ \\
S1, roll 5° & \cellcolor{mejor}$-34.4$ & \cellcolor{mejor}$-32.9$/$-34.5$ & \cellcolor{mejor}$-31.0$/$-33.3$ & \cellcolor{regular}$-27.3$/$-29.8$ & \cellcolor{peor}$-20.1$/$-23.0$ & & --- \\
S2 ($\theta_{pB}=68.66^{\circ}$) & \cellcolor{mejor}$-30.7$ & & & & & & $-12.6$ \\
S3 ($\theta_{pB}=74.05^{\circ}$) & \cellcolor{mejor}$-26.5$ & & & & & & $-8.7$ \\
\bottomrule
\end{tabular}}
\end{table}

\begin{itemize}[itemsep=2pt]
  \item En el suelo sin pérdidas el mínimo es un cero exacto, pero \textbf{desaparece
  con 0.5°} de error ($-190\to-42$\,dB): es el ``mínimo extremadamente estrecho'' que
  §2.1 advierte que no basta para justificar una trayectoria.
  \item Con pérdidas el mínimo es \textbf{poco profundo pero ancho}: en S1 sólo
  $-34.5$\,dB, y se pierden 0.7, 2.4, 6.0 y 13.0\,dB con errores de 0.5°, 1°, 2° y 5°.
  Respecto de la incidencia base (53.4°), la supresión disponible es de 17.0\,dB.
  El mínimo empeora con la humedad ($-30.7$ y $-26.5$\,dB en S2 y S3) y se desplaza
  a 68.7° y 74.0°.
  \item El roll de 1°--5° es irrelevante en S1 (variación $\le1.5$\,dB); en S0 sólo
  desplaza el cero $\approx0.3^{\circ}$ (modelo co-polar, sin término cruzado).
\end{itemize}

\paragraph{Brewster biestático del clutter difuso.} La ec.~\eqref{eq:spm} muestra
que, para el retorno difuso, el pseudo-Brewster no es un punto sino una \emph{curva} en el
espacio (incidencia, separación): para $\theta<\theta_{pB}$ existe una separación
$\Delta\phi^\ast(\theta)<180^{\circ}$ donde el retorno difuso VV de primer orden se anula.
La curva termina en el Brewster clásico ($\Delta\phi=180^{\circ}$, $\theta=\theta_{pB}$) y
nunca pasa por la mono ($\Delta\phi=0$). En HH el nulo está en $\Delta\phi=90^{\circ}$ para
cualquier incidencia y es puramente geométrico. Comparar TM con TE en la misma geometría
separa ambos efectos (§5.4).

\begin{figure}[H]
\centering
\includegraphics[width=\textwidth]{C_sigma0_mapa_suelos.pdf}
\caption{$\sigma^0$ biestático SPM con el Rx a la misma elevación que el Tx, para los
cuatro suelos ($s=2$\,cm). Arriba VV (TM), abajo HH (TE). Curva roja: mínimo de VV en cada
incidencia, $\cos\phi_s=\varepsilon\sin^2\theta/(\varepsilon-\sin^2\theta)$; vertical roja:
nulo de HH. Discontinua: $\theta_{pB}$ de cada suelo; banda clara: incidencias de la hélice base.}
\label{fig:C_sigma0_suelos}
\end{figure}

El comportamiento es el mismo en los cuatro suelos (Fig.~\ref{fig:C_sigma0_suelos}):
\begin{itemize}[itemsep=1pt]
  \item \textbf{VV:} en la incidencia media de la hélice (53.4°) el mínimo está en
  $\Delta\phi^\ast=140^{\circ}$ (S0, S1), 136° (S2) y 133° (S3). Está 34--39\,dB por debajo de la mono
  en los suelos con pérdidas (en S0, 74\,dB) y se mantiene a menos de 10\,dB del mínimo en una
  franja de 4--8° de $\Delta\phi$ (0.1° en S0). Al subir la humedad se desplaza poco y se estrecha.
  \item \textbf{HH:} nulo exacto en $\Delta\phi=90^{\circ}$ para todos los suelos e incidencias, porque
  $\alpha_{hh}\propto\cos\phi_s$ no depende de $\varepsilon_r$: es el desacople geométrico entre los
  vectores TE del Tx y del Rx, ortogonales a 90°.
\end{itemize}
La sección siguiente muestra cómo se reflejan estos nulos en el SCR del punto escalar, y
la Sección~\ref{sec:C2} comprueba si sobreviven con blancos reales.

\section{SCR en la imagen}

\begin{figure}[H]
\centering
\includegraphics[width=\textwidth]{C_scr_incidencia_rugosa.pdf}
\caption{SCR en el vóxel del blanco frente a la incidencia central (S1, $s=2$\,cm),
familias F1 y F2, TM y TE. Línea roja: pseudo-Brewster.}
\label{fig:C_scr}
\end{figure}

\begin{figure}[H]
\centering
\includegraphics[width=\textwidth]{C_scr_error_angular.pdf}
\caption{SCR frente al error de incidencia respecto de $\theta_{pB}$ (F2, S1, $s=2$\,cm).}
\label{fig:C_err}
\end{figure}

\begin{table}[H]
\centering\small
\caption{SCR en imagen [dB] en el vóxel del blanco: S1, $s=2$\,cm, superficie congelada (3 realizaciones). \colorbox{mejor}{\textbf{verde}}: mejor canal de la fila; \colorbox{peor}{rojo}: SCR $<0$ (la superficie domina). $\leftarrow$: fila en el pseudo-Brewster. CSV: \texttt{expC/expC\_scr\_imagen.csv}.}
\label{tab:C_scr}
\begin{tabular}{@{}llrrrrrr@{}}
\toprule
Pol. & Familia, $\theta_c$ & Mono & Bi 90° & Bi 135° & Ref1 180° & Mono+bi 90° & Mono+bi 135° \\
\midrule
TM & F1, 35° & \cellcolor{peor}$-$0.5 & 3.3 & 0.8 & \cellcolor{peor}$-$59.6 & 2.8 & \cellcolor{mejor}\textbf{5.8} \\
 & F1, 45° & \cellcolor{peor}$-$0.2 & 1.3 & \cellcolor{mejor}\textbf{13.0} & \cellcolor{peor}$-$58.2 & 4.4 & 5.6 \\
 & F1, 53.4° & 3.1 & 3.1 & \cellcolor{mejor}\textbf{9.8} & \cellcolor{peor}$-$55.3 & 4.3 & 9.5 \\
 & \textbf{F1, $\theta_{pB}$=63.45°} $\leftarrow$ & 8.0 & 6.0 & 7.4 & \cellcolor{peor}$-$41.5 & 9.4 & \cellcolor{mejor}\textbf{12.0} \\
 & F1, 75° & 14.8 & 11.6 & 12.4 & \cellcolor{peor}$-$66.3 & 15.6 & \cellcolor{mejor}\textbf{18.0} \\
\cmidrule(l){2-8}
 & F2, 35° & \cellcolor{peor}$-$4.7 & \cellcolor{peor}$-$1.6 & \cellcolor{peor}$-$13.4 & \cellcolor{peor}$-$61.2 & \cellcolor{mejor}\textbf{$-$0.2} & \cellcolor{peor}$-$8.3 \\
 & F2, 45° & 0.5 & \cellcolor{peor}$-$0.9 & 3.4 & \cellcolor{peor}$-$60.1 & 1.6 & \cellcolor{mejor}\textbf{5.6} \\
 & F2, 53.4° & 0.2 & 3.0 & \cellcolor{mejor}\textbf{6.9} & \cellcolor{peor}$-$57.3 & 2.9 & 4.2 \\
 & \textbf{F2, $\theta_{pB}$=63.45°} $\leftarrow$ & 2.5 & 3.5 & 5.1 & \cellcolor{peor}$-$41.5 & \cellcolor{mejor}\textbf{7.8} & 7.3 \\
 & F2, 75° & 8.4 & 5.7 & 8.3 & \cellcolor{peor}$-$66.0 & \cellcolor{mejor}\textbf{12.6} & 12.1 \\
\midrule
TE & F1, 35° & 1.4 & \cellcolor{mejor}\textbf{14.4} & 1.5 & \cellcolor{peor}$-$64.4 & 6.7 & 4.2 \\
 & F1, 45° & 2.4 & \cellcolor{mejor}\textbf{18.6} & \cellcolor{peor}$-$4.9 & \cellcolor{peor}$-$66.9 & 8.8 & 0.7 \\
 & F1, 53.4° & 6.2 & \cellcolor{mejor}\textbf{23.8} & \cellcolor{peor}$-$2.3 & \cellcolor{peor}$-$69.8 & 11.6 & 4.6 \\
 & \textbf{F1, $\theta_{pB}$=63.45°} $\leftarrow$ & 11.5 & \cellcolor{mejor}\textbf{28.9} & \cellcolor{peor}$-$1.3 & \cellcolor{peor}$-$74.4 & 17.1 & 4.7 \\
 & F1, 75° & 18.5 & \cellcolor{mejor}\textbf{35.0} & 6.9 & \cellcolor{peor}$-$81.3 & 24.6 & 13.6 \\
\cmidrule(l){2-8}
 & F2, 35° & \cellcolor{peor}$-$2.6 & \cellcolor{mejor}\textbf{11.7} & \cellcolor{peor}$-$18.1 & \cellcolor{peor}$-$65.8 & 3.6 & \cellcolor{peor}$-$12.0 \\
 & F2, 45° & 3.5 & \cellcolor{mejor}\textbf{18.5} & \cellcolor{peor}$-$9.1 & \cellcolor{peor}$-$68.6 & 9.0 & \cellcolor{peor}$-$3.3 \\
 & F2, 53.4° & 3.8 & \cellcolor{mejor}\textbf{24.6} & \cellcolor{peor}$-$3.9 & \cellcolor{peor}$-$71.5 & 9.5 & 0.2 \\
 & \textbf{F2, $\theta_{pB}$=63.45°} $\leftarrow$ & 6.7 & \cellcolor{mejor}\textbf{28.7} & \cellcolor{peor}$-$1.9 & \cellcolor{peor}$-$75.5 & 12.9 & 4.2 \\
 & F2, 75° & 13.1 & \cellcolor{mejor}\textbf{36.0} & 4.4 & \cellcolor{peor}$-$81.6 & 19.4 & 10.2 \\
\bottomrule
\end{tabular}
\end{table}


\paragraph{Lectura de la Tabla~\ref{tab:C_scr}.} Cuanto más alto el SCR, mejor.
\begin{enumerate}[itemsep=3pt]
  \item \textbf{El pseudo-Brewster sólo se nota en Ref1, y Ref1 es inservible.} Es la única
  columna con un pico en $\theta_{pB}$: en TM pasa de $-55.3$ a $-41.5$\,dB, y en TE no hay pico
  (control). Aun así queda más de 40\,dB por debajo de cualquier otro canal (toda la columna en rojo).
  Con espectro rectangular la conclusión es la misma ($-34.4$\,dB en $\theta_{pB}$).
  \item \textbf{En los demás canales, más incidencia es mejor, sin pico en $\theta_{pB}$.} El SCR
  crece de forma casi monótona de 35° a 75°, igual en TM que en TE. Por ejemplo, mono F1:
  3.1\,$\to$\,8.0\,$\to$\,14.8\,dB en TM y 6.2\,$\to$\,11.5\,$\to$\,18.5\,dB en TE, en
  53.4°, 63.45° y 75°. La causa es que la difusa cae con la incidencia, por la longitud de
  correlación $\ell$ (factor $W(K)$), no Brewster. El máximo de casi todas las columnas está en
  75°, la mayor incidencia probada.
  \item \textbf{TE supera a TM.} En mono TE gana 2--5\,dB en todas las incidencias,
  al contrario de lo que se esperaría si Brewster (TM) importara: en retrodispersión la difusa
  VV es mayor que la HH.
  \item \textbf{Mejor canal con el punto escalar.} En TE, el bi a 90° es el mejor en las 10 filas
  (12--36\,dB según la incidencia), por el nulo de HH a 90°. \emph{Este resultado no se
  sostiene con blancos reales} (Sección~\ref{sec:C2}). En TM, un canal a 135° (bi o mono+bi) es el mejor en 7 de 10 filas, con 3--13\,dB
  sobre la mono.
\end{enumerate}

\begin{table}[H]
\centering\small
\caption{F1 (hélice) frente a F2 (cono): $\mathrm{SCR}_{F1}-\mathrm{SCR}_{F2}$ en la misma incidencia central, polarización y canal (10 pares por canal: 5 incidencias $\times$ TM/TE, Tabla~\ref{tab:C_scr}). \colorbox{mejor}{verde}: F1 gana.}
\label{tab:C_f1f2}
\begin{tabular}{@{}lrrrc@{}}
\toprule
Canal & media [dB] & mín. [dB] & máx. [dB] & F1 gana \\
\midrule
Mono & \cellcolor{mejor}$+$3.4 & $-$1.0 & 6.4 & 8/10 \\
Bi 90° & \cellcolor{mejor}$+$1.7 & $-$1.0 & 5.9 & 8/10 \\
Bi 135° & \cellcolor{mejor}$+$6.2 & 0.7 & 19.5 & 10/10 \\
Ref1 180° & \cellcolor{mejor}$+$1.2 & $-$0.2 & 2.1 & 9/10 \\
Mono+bi 90° & \cellcolor{mejor}$+$2.6 & $-$0.2 & 5.3 & 9/10 \\
Mono+bi 135° & \cellcolor{mejor}$+$5.9 & 0.0 & 16.2 & 10/10 \\
\midrule
Todos & $+$3.5 & & & 54/60 \\
\bottomrule
\end{tabular}
\end{table}


\paragraph{F1 frente a F2.} La hélice (F1) supera al cono (F2) en 54 de las 60 comparaciones
de igual incidencia, polarización y canal, con 3.5\,dB de SCR de media
(Tabla~\ref{tab:C_f1f2}). El eco del blanco es idéntico en las dos familias (diferencia
$\le0.2$\,dB): lo que cambia es cuánto clutter de la superficie llega al vóxel del blanco. La
superficie está 2\,m por encima del blanco, y F2 tiene peor resolución vertical
($\delta_z$ 2.34 frente a 2.08\,m) y peor ISLR (Tabla~\ref{tab:C_costo}), de modo que filtra
más energía de la superficie hacia el blanco. F2, la órbita que impone Brewster, es por tanto
peor también en SCR, la métrica para la que se propuso.

\begin{table}[H]
\centering\small
\caption{Robustez del SCR en imagen [dB] frente al error de incidencia $\delta$ alrededor de $\theta_{pB}$ (TM, $s=2$\,cm). Pares $a$\,/\,$b$: $\delta<0$\,/\,$\delta>0$. \colorbox{mejor}{verde}: mejor canal de la fila; \colorbox{peor}{rojo}: SCR $<0$.}
\label{tab:C_err_scr}
\begin{tabular}{@{}llrrrr@{}}
\toprule
Suelo, familia & $|\delta|$ & Mono & Bi 135° & Ref1 180° & Mono+bi 135° \\
\midrule
S1, F1 & 0° & 8.0 & 7.4 & \cellcolor{peor}$-$41.5 & \cellcolor{mejor}\textbf{12.0} \\
 & 0.5° & 7.3 / 6.8 & 10.4 / 8.5 & \cellcolor{peor}$-$40.4 / $-$43.4 & \cellcolor{mejor}\textbf{11.9 / 10.5} \\
 & 1° & 6.5 / 10.6 & 8.4 / 9.6 & \cellcolor{peor}$-$40.9 / $-$45.4 & \cellcolor{mejor}\textbf{11.3 / 12.9} \\
 & 2° & 6.6 / 9.6 & 9.4 / 12.3 & \cellcolor{peor}$-$43.9 / $-$49.0 & \cellcolor{mejor}\textbf{11.3 / 15.9} \\
 & 5° & 2.8 / 13.4 & 5.2 / 9.2 & \cellcolor{peor}$-$50.6 / $-$56.4 & \cellcolor{mejor}\textbf{7.0 / 13.6} \\
\midrule
S1, F2 & 0° & 2.5 & 5.1 & \cellcolor{peor}$-$41.5 & \cellcolor{mejor}\textbf{7.3} \\
 & 0.5° & 4.0 / 2.5 & 4.6 / 6.8 & \cellcolor{peor}$-$42.2 / $-$42.4 & \cellcolor{mejor}\textbf{8.3 / 6.4} \\
 & 1° & 4.7 / 5.0 & 3.2 / 7.9 & \cellcolor{peor}$-$43.8 / $-$44.2 & \cellcolor{mejor}\textbf{7.8 / 8.7} \\
 & 2° & 2.5 / 5.4 & 7.2 / 7.9 & \cellcolor{peor}$-$47.1 / $-$48.2 & \cellcolor{mejor}\textbf{8.5 / 9.3} \\
 & 5° & 4.8 / 6.0 & 5.1 / 5.3 & \cellcolor{peor}$-$53.0 / $-$56.2 & \cellcolor{mejor}\textbf{9.6 / 8.7} \\
\midrule
S0, F2 & 0° & 39.0 & 35.9 & \cellcolor{mejor}\textbf{61.2} & 43.1 \\
 & 2° & 37.8 / 41.8 & 38.9 / 40.8 & \cellcolor{peor}$-$11.3 / $-$12.4 & \cellcolor{mejor}\textbf{41.7 / 46.6} \\
\bottomrule
\end{tabular}
\end{table}


\paragraph{Robustez al error de incidencia (Tabla~\ref{tab:C_err_scr}).} En S1, con F2, errores
de hasta $\pm5^{\circ}$ cambian el SCR de mono, bi y fusión en unos $\pm3$\,dB sin tendencia (del
orden de la dispersión entre realizaciones). Con F1 la variación llega a $\pm5$\,dB, pero sigue
la tendencia general con la incidencia ($+\delta$ mejora, $-\delta$ empeora), no un efecto de
Brewster. Sólo Ref1 cae al alejarse de $\theta_{pB}$, hasta 15\,dB con $\pm5^{\circ}$. Es decir,
fuera de Ref1 el SCR es robusto al error angular \emph{porque no depende de Brewster}. En el suelo sin pérdidas
(S0), Ref1 llega a 61\,dB en $\theta_{pB}$ exacto y cae a $-12$\,dB con $\pm2^{\circ}$: es el
mínimo extremadamente estrecho que §2.1 descarta.

\paragraph{Otros factores (una línea cada uno).}
\begin{itemize}[itemsep=1pt]
  \item \textbf{Rugosidad:} el SCR de mono y bi intermedio cae $\approx6$\,dB cada vez que $s$ se
  duplica (difusa $\propto s^2$); el de Ref1 casi no cambia (domina la especular).
  \item \textbf{Suelos húmedos} (S2, S3): SCR negativo en todas las geometrías ($-14$ a $-28$\,dB
  en S2 y $-58$ a $-84$\,dB en S3), porque domina la atenuación del blanco a 2\,m (13.6 y
  27.7\,dB/m por trayecto, frente a 7.7\,dB/m en S1). El pseudo-Brewster no lo remedia.
  \item \textbf{Atitud:} un roll de 0--5° no cambia el SCR en S1 ($\le0.2$\,dB).
\end{itemize}

\section{Costo en imagen de la restricción de incidencia}

\begin{figure}[H]
\centering
\includegraphics[width=\textwidth]{C_costo_imagen.pdf}
\caption{Resolución, volumen y PSLR frente a la incidencia central para F1 (continuo)
y F2 (discontinuo), mono y mono+bi 135°.}
\label{fig:C_costo}
\end{figure}

\begin{table}[H]
\centering\small
\caption{Hélice base (F1 en 53.4°) frente a la restricción de pseudo-Brewster (F2 en
$\theta_{pB}$), con ponderación Hann (\texttt{expC\_psf\_vs\_incidencia.csv},
\texttt{expC\_scr\_imagen.csv}). En cada canal, \colorbox{mejor}{verde} marca la trayectoria que
gana cada métrica. Menor es mejor salvo en contraste y SCR.}
\label{tab:C_costo}
\begin{tabular}{@{}llrrrrrr@{}}
\toprule
Trayectoria & Canal & $\delta_x$ [m] & $\delta_z$ [m] & $V_{-3}$ [l] & ISLR [dB] & Contraste [dB] & SCR [dB] \\
\midrule
F1 base (53.4°) & mono & 0.157 & \cellcolor{mejor}2.08 & 29.0 & \cellcolor{mejor}6.5 & \cellcolor{mejor}19.6 & \cellcolor{mejor}3.1 \\
F2 cono $\theta_{pB}$ & mono & \cellcolor{mejor}0.140 & 2.34 & \cellcolor{mejor}25.0 & 8.3 & 18.4 & 2.5 \\
\midrule
F1 base (53.4°) & mono+bi 135° & 0.213 & \cellcolor{mejor}2.08 & 52.0 & \cellcolor{mejor}2.4 & \cellcolor{mejor}20.4 & \cellcolor{mejor}9.5 \\
F2 cono $\theta_{pB}$ & mono+bi 135° & \cellcolor{mejor}0.191 & 2.34 & \cellcolor{mejor}46.0 & 3.6 & 19.9 & 7.3 \\
\bottomrule
\end{tabular}
\end{table}

\textbf{Costo (pregunta 7).} Imponer la órbita de pseudo-Brewster (F2 en $\theta_{pB}$) sólo gana en
horizontal: $\delta_x$ mejora un 11\,\% (mayor $\sin\theta$) y el volumen del lóbulo un
12--14\,\%. Pierde en todo lo demás: $\delta_z$ empeora un 12.5\,\% (se pierde la diversidad de
elevación de la hélice), el ISLR sube 1.2--1.8\,dB, el contraste baja 0.5--1.2\,dB y el SCR baja
0.6--2.2\,dB. Además obliga a bajar la plataforma a $\approx85$\,m.

\section{Exp.~C2: ¿sobreviven los nulos con blancos reales?}
\label{sec:C2}

\paragraph{Motivación.} El SCR de la Tabla~\ref{tab:C_scr} usa un punto escalar ($S=1$), que no
tiene polarización: sólo la superficie siente los nulos de VV y HH. Un blanco real también
dispersa con polarización, y el mismo desacople geométrico puede anular su eco. El Exp.~C2
separa ambos efectos.

\paragraph{Método.} Tx en la hélice base (53.4°), suelo S1, $s=2$\,cm, $\ell=20$\,cm,
ponderación Hann y 3 realizaciones de la superficie congelada. Rx en la misma hélice con
$\Delta\phi=0^{\circ}{:}15^{\circ}{:}180^{\circ}$ y 140°, TM y TE co-polares. Cuatro blancos en $(0,0,-2)$\,m:
punto escalar, esfera PEC de $r=0.5$\,m, cubo PEC de 1\,m en orientación nominal y el mismo cubo
en 24 orientaciones SO(3). Métricas: SCR en el vóxel del blanco, como en el Exp.~C, y PSCR,
el pico del blanco bajo la superficie dividido por el pico de la superficie en toda la rejilla,
para el canal bi aislado y para la fusión con el mismo presupuesto (N). Script
\texttt{exp\_C2\_nulos.py}; CSV \texttt{expC2/expC2\_scr\_vs\_dphi.csv}. La rejilla y las
realizaciones difieren de las del Exp.~C, de modo que los valores absolutos del punto difieren
hasta 3\,dB de la Tabla~\ref{tab:C_scr}; las comparaciones se hacen dentro de esta tabla.

\begin{figure}[H]
\centering
\includegraphics[width=\textwidth]{C2_nulos_scr.pdf}
\caption{Exp.~C2 frente a $\Delta\phi$ (hélice base, S1, $s=2$\,cm). (a, d) Potencia en la
imagen del canal bi respecto de la mono: superficie en el vóxel del blanco (rojo) y eco de cada
blanco. (b, e) Ganancia de SCR en el vóxel respecto de la mono. (c, f) Ganancia de PSCR (pico
del blanco / pico de la superficie). Continua: bi aislado; discontinua: fusión mono+bi (N).
Líneas grises en 90° (nulo HH) y 140° (nulo VV).}
\label{fig:C2}
\end{figure}

\begin{table}[H]
\centering\small
\caption{Exp.~C2: SCR en el vóxel del blanco [dB]. ``Superficie'': potencia de la superficie
en el vóxel respecto de la mono. \colorbox{mejor}{verde}: mejora de más de 3\,dB sobre la mono
de la misma polarización; \colorbox{peor}{rojo}: empeora más de 3\,dB.}
\label{tab:C2}
\begin{tabular}{@{}llrrrrr@{}}
\toprule
Pol. & Canal & Superficie & Punto & Esfera & Cubo nominal & Cubo SO(3) \\
\midrule
TM & mono & 0.0 & 0.7 & $-0.9$ & $-9.5$ & $-12.0$ \\
 & bi 90° & $-2.0$ & 2.6 & \cellcolor{peor}$-13.3$ & \cellcolor{peor}$-25.3$ & \cellcolor{peor}$-27.9$ \\
 & bi 120° & $-4.6$ & \cellcolor{mejor}5.2 & \cellcolor{peor}$-6.1$ & \cellcolor{peor}$-19.1$ & \cellcolor{mejor}$-6.2$ \\
 & bi 135° & $-12.5$ & \cellcolor{mejor}13.1 & \cellcolor{mejor}7.5 & \cellcolor{mejor}14.7 & \cellcolor{mejor}6.2 \\
 & bi 140° & $-13.6$ & \cellcolor{mejor}14.2 & \cellcolor{mejor}\textbf{9.8} & \cellcolor{mejor}\textbf{20.1} & \cellcolor{mejor}\textbf{8.3} \\
 & bi 150° & $-1.2$ & 1.8 & $-0.9$ & \cellcolor{mejor}13.7 & \cellcolor{mejor}$-2.7$ \\
 & mono+bi 140° & $-5.7$ & \cellcolor{mejor}6.4 & $-0.9$ & \cellcolor{mejor}4.7 & \cellcolor{mejor}$-4.2$ \\
\midrule
TE & mono & 0.0 & 3.9 & 2.3 & $-6.4$ & $-8.8$ \\
 & bi 90° & $-17.8$ & \cellcolor{mejor}\textbf{21.7} & \cellcolor{peor}$-4.8$ & $-6.2$ & $-9.0$ \\
 & mono+bi 90° & $-5.5$ & \cellcolor{mejor}9.4 & 2.2 & $-7.4$ & $-9.3$ \\
\bottomrule
\end{tabular}
\end{table}

\paragraph{Resultados e interpretación.}
\begin{enumerate}[itemsep=3pt]
  \item \textbf{El nulo HH en TE a 90° es un artefacto del punto escalar.} La superficie cae
  17.8\,dB, pero el eco de la esfera cae 25\,dB y el del cubo SO(3) 18\,dB. El SCR de la esfera
  empeora 7.1\,dB respecto de la mono y el del cubo SO(3) no cambia. Es el mismo mecanismo que el
  de la Sección~\ref{sec:D3}: a 90° los vectores TE del Tx y del Rx son ortogonales, y eso
  desacopla tanto la superficie como el blanco. Sólo el punto, sin polarización, gana 17.8\,dB.
  \item \textbf{El nulo VV en TM a 135°--140° sí sobrevive.} La superficie cae 12.5 a 13.6\,dB,
  mientras que el eco de los blancos apenas cae (esfera, $-4.1$ y $-2.9$\,dB) o incluso sube
  (cubo SO(3), $+5.8$ y $+6.7$\,dB, por la especular de sus caras). El SCR del canal bi aislado
  mejora respecto de la mono 8.4 a 10.7\,dB en la esfera, 24 a 30\,dB en el cubo nominal y 18 a
  20\,dB en el cubo con orientación aleatoria. En PSCR la mejora es menor pero del mismo signo:
  $+2.5$ a $+3.7$\,dB en la esfera y $+12$ a $+14$\,dB en el cubo SO(3).
  \item \textbf{La ventana es estrecha.} A 120° y a 150° la esfera ya no mejora ($-6.1$ y
  $-0.9$\,dB) y desde 150° la especular de nadir domina el PSCR. El Rx debe mantenerse a
  $\approx\pm5^{\circ}$ de 140°, coherente con la anchura de 4--8° del nulo VV
  (Fig.~\ref{fig:C_sigma0_suelos}).
  \item \textbf{La fusión de igual peso desperdicia el canal limpio.} La imagen fusionada es la
  media de las dos, y el clutter de la mono la domina: su superficie sólo baja 5.7\,dB, casi el
  límite de 6\,dB de una media de dos canales. A 140° la fusión deja la esfera igual que la mono y
  queda 10 a 15\,dB por debajo del bi aislado. Para aprovechar el nulo hay que usar el canal bi
  solo o una fusión ponderada por el clutter de cada canal.
  \item \textbf{El precio es la resolución.} A 135° el canal bi tiene $\delta_x=0.41$\,m frente a
  0.157\,m en la mono (Tabla~\ref{tab:expB}), con la misma $\delta_z$. Es un modo de detección
  en escenarios limitados por clutter, no un modo de imagen fina.
\end{enumerate}

\conclexp{C2}{Con blancos reales, la única geometría que suprime el clutter de forma útil es
TM con el Rx a 135°--140° usado como canal aislado: el SCR mejora 8 a 20\,dB sobre la mono, a
cambio de una resolución horizontal 2.6 veces peor. El nulo HH en TE a 90° es un artefacto del
punto escalar.}

\section{Conclusiones del Exp.~C: ¿qué es mejor?}

\begin{table}[H]
\centering\small
\caption{Respuesta directa a las preguntas del Exp.~C (S1, $\varepsilon_r=3.99-j0.40$, $s=2$\,cm).
\colorbox{mejor}{verde}: la opción recomendada.}
\label{tab:C_concl}
\begin{tabular}{@{}L{3.6cm}L{3.2cm}L{3.4cm}L{4.4cm}@{}}
\toprule
\textbf{Pregunta} & \textbf{Mejor opción} & \textbf{Margen} & \textbf{Por qué} \\
\midrule
¿Volar restringido al pseudo-Brewster? & \cellcolor{mejor}No & F2 en $\theta_{pB}$ pierde 0.6--2.2\,dB de SCR (mono y mono+bi 135°) frente a la hélice base &
Sólo suprime la especular, que sólo pesa en Ref1 \\
¿Hélice (F1) o cono (F2)? & \cellcolor{mejor}F1 (hélice) & $+3.5$\,dB de SCR medio; gana 54/60 &
Mejor $\delta_z$ e ISLR: menos superficie en el vóxel del blanco \\
¿Qué incidencia? & \cellcolor{mejor}La mayor viable (75° fue la mejor probada) & Mono TM: $+11.7$\,dB de 53.4° a 75° &
La difusa cae con $\theta$ ($\ell$); igual en TM y TE \\
¿Qué canal en TM? & \cellcolor{mejor}Bi aislado a 135°--140° & 8--20\,dB de SCR sobre la mono con esfera y cubo SO(3) (Exp.~C2) & Nulo VV biestático; el eco del blanco no cae (SPM; hay que validarlo) \\
¿TM o TE? & \cellcolor{mejor}TM & TE a 90° sólo gana con el punto escalar; con la esfera pierde 7\,dB & El desacople a 90° anula también el eco del blanco \\
¿Fusión de igual peso? & \cellcolor{peor}No para clutter & Queda 10--15\,dB por debajo del bi aislado a 140° & El clutter de la mono domina la media \\
¿Ref1 (180°)? & \cellcolor{peor}Nunca & $-41$ a $-81$\,dB & La especular cae sobre el blanco \\
¿Es robusto al error angular? & Sí, salvo Ref1 & F2: $\pm3$\,dB con $\pm5^{\circ}$ & Porque no depende de Brewster \\
\bottomrule
\end{tabular}
\end{table}

\paragraph{En una frase.} El pseudo-Brewster no es un criterio útil de trayectoria: la órbita
que lo impone (F2) es peor que la hélice en SCR, resolución vertical y lóbulos, y sólo mejora
Ref1, que sigue inservible. Lo que sí sube el SCR es la geometría, no Brewster: mayor incidencia y, sobre todo,
el Rx a 135°--140° en TM usado como canal aislado, que mejora 8 a 20\,dB el SCR de blancos reales.

\paragraph{Cautelas.} (i) La incidencia de 75° se evaluó sin restricciones operativas: el
Tx volaría a $\approx49$\,m de altura, con más alcance y otra cobertura del haz. (ii) El nulo
de VV es de SPM de primer orden y debe verificarse con un modelo de orden superior (IEM/AIEM),
con polarización cruzada y con medidas. (iii) El Exp.~C2 usa 3 realizaciones de la superficie y
24 orientaciones del cubo; las ganancias de 8 a 20\,dB superan con margen la dispersión de
$\pm3$\,dB.
